\documentclass[../../main.tex]{subfiles}
\begin{document}

{\bf [解答]}

% まずは$m$個のビーカーの様子
\begin{figure}[htb]
  \centering
  \begin{tikzpicture}[scale=1.0, >=stealth, every node/.style={font=\small}]

    % --- ビーカー描画マクロ ---
    % 引数: #1=中心x, #2=水深比率(0~1), #3=ビーカー名, #4=水量ラベル, #5=ハイライト色
    \newcommand{\drawBeaker}[5]{
      % 水の塗りつぶし
      \fill[#5!25] ({#1-0.5}, 0) rectangle ({#1+0.5}, {#2*1.8});
      % 水面線
      \draw[#5!80!black, thick] ({#1-0.5}, {#2*1.8}) -- ({#1+0.5}, {#2*1.8});
      % ビーカー枠線（上側は開いたU字形）
      \draw[thick, gray!75!black] ({#1-0.5}, 2.0) -- ({#1-0.5}, 0) -- ({#1+0.5}, 0) -- ({#1+0.5}, 2.0);
      % ラベル
      \node[below=2pt, font=\scriptsize\bfseries] at ({#1}, 0) {#3};
      \node[above=2pt, font=\scriptsize, #5!90!black] at ({#1}, {#2*1.8}) {#4};
    }

    % --- 1. ビーカーの並び ---
    % ビーカー 1
    \drawBeaker{0.8}{0.35}{ビーカー $1$}{$v_1$}{blue}
    % ビーカー 2
    \drawBeaker{2.3}{0.20}{ビーカー $2$}{$v_2$}{blue}
    % 点々（省略）
    \node[font=\large] at (3.5, 0.9) {$\dots$};
    % ビーカー i
    \drawBeaker{4.7}{0.45}{ビーカー $i$}{$v_i$}{blue}
    % 点々（省略）
    \node[font=\large] at (5.9, 0.9) {$\dots$};
    % ビーカー X（最大水量 x を持つ特別なビーカー: 赤系ハイライト）
    \drawBeaker{7.3}{0.75}{ビーカー $X$}{$x$}{red}
    % 点々（省略）
    \node[font=\large] at (8.5, 0.9) {$\dots$};
    % ビーカー m
    \drawBeaker{9.7}{0.30}{ビーカー $m$}{$v_m$}{blue}

    % --- 2. 全体括れ線と総水量の表示 ---
    \draw[decorate, decoration={brace, amplitude=6pt}, thick, gray!70!black]
    (0.2, 2.3) -- (10.3, 2.3)
    node[midway, above=7pt, font=\footnotesize\bfseries] {全 $m$ 個のビーカー（総水量 $S = 1\text{ L}$）};

  \end{tikzpicture}
  \caption{初期状態における $m$ 個のビーカーと総水量 $S = 1\text{ L}$}
  \label{fig:beakers_initial_state}
\end{figure}

$x$ [L] の水が入っているビーカーを $X$ とする．
ビーカーに入っている水の全質量を $S$ [L]とする．題意より
\begin{align}
  S = 1 \label{eq:1}
\end{align}
である．

% 操作の様子
\begin{figure}[htb]
  \centering
  \begin{tikzpicture}[scale=1.1, >=stealth, every node/.style={font=\small}]

    % --- ビーカー描画マクロ ---
    \newcommand{\drawMiniBeaker}[5]{
      \fill[#5!25] ({#1-0.45}, 0) rectangle ({#1+0.45}, {#2*1.5});
      \draw[#5!80!black, thick] ({#1-0.45}, {#2*1.5}) -- ({#1+0.45}, {#2*1.5});
      \draw[thick, gray!75!black] ({#1-0.45}, 1.7) -- ({#1-0.45}, 0) -- ({#1+0.45}, 0) -- ({#1+0.45}, 1.7);
      \node[below=2pt, font=\scriptsize] at ({#1}, 0) {#3};
      \node[above=1pt, font=\tiny, #5!90!black] at ({#1}, {#2*1.5}) {#4};
    }

    % ============================================================
    % 左側: 操作前（最少の2つを選択）
    % ============================================================
    \begin{scope}[xshift=0cm]
      \node[font=\footnotesize\bfseries, blue!85!black] at (1.6, 2.3) {【操作前】水量の少ない2つを選択};

      % ビーカー A (最少: 例 0.2L)
      \drawMiniBeaker{0.6}{0.20}{(a) 最少}{$v_{\min1}$}{orange}
      % ビーカー B (2番目に少: 例 0.35L)
      \drawMiniBeaker{1.8}{0.35}{(b) 2番目}{$v_{\min2}$}{orange}
      % その他のビーカー
      \drawMiniBeaker{3.0}{0.65}{他}{}{blue}

      % 注ぎ入れる矢印
      \draw[->, line width=1.5pt, orange!90!black] (0.6, 1.8) to[out=60, in=120] (1.8, 1.8);
      \node[font=\tiny, orange!95!black, above=1pt] at (1.2, 2.0) {合流};
    \end{scope}

    % 中央の操作遷移矢印
    \draw[->, line width=2.0pt, gray!60] (3.8, 0.9) -- (4.8, 0.9)
    node[midway, above=2pt, font=\scriptsize\bfseries, gray!80!black] {1回操作};

    % ============================================================
    % 右側: 操作後（合流してビーカー数が1つ減少）
    % ============================================================
    \begin{scope}[xshift=5.3cm]
      \node[font=\footnotesize\bfseries, teal!85!black] at (1.6, 2.3) {【操作後】ビーカー数が1個減少};

      % 空になって除外されるビーカー
      \draw[dashed, gray!50, thick] (0.15, 1.7) -- (0.15, 0) -- (1.05, 0) -- (1.05, 1.7);
      \node[below=2pt, font=\scriptsize, gray!60] at (0.6, 0) {空（除外）};
      \node[font=\tiny, gray!50] at (0.6, 0.8) {除外};

      % 水が合わさったビーカー (0.2 + 0.35 = 0.55L)
      \drawMiniBeaker{1.8}{0.55}{合流後}{$v_{\min1} + v_{\min2}$}{teal}
      % 残りのビーカー
      \drawMiniBeaker{3.0}{0.65}{他}{}{blue}
    \end{scope}

  \end{tikzpicture}
  \caption{水量の少ない2つのビーカーを合流させる操作}
  \label{fig:beaker_operation_step}
\end{figure}

\paragraph{(1)}

% 残り3つで，t,s,xの水が入っている様子
\begin{figure}[htb]
  \centering
  \begin{tikzpicture}[scale=1.3, >=stealth, every node/.style={font=\small}]

    % --- ビーカー描画マクロ ---
    \newcommand{\drawFinalBeaker}[5]{
      \fill[#5!25] ({#1-0.6}, 0) rectangle ({#1+0.6}, {#2*2.2});
      \draw[#5!80!black, thick] ({#1-0.6}, {#2*2.2}) -- ({#1+0.6}, {#2*2.2});
      \draw[very thick, gray!75!black] ({#1-0.6}, 2.3) -- ({#1-0.6}, 0) -- ({#1+0.6}, 0) -- ({#1+0.6}, 2.3);
      \node[below=3pt, font=\footnotesize\bfseries] at ({#1}, 0) {#3};
      \node[above=2pt, font=\footnotesize\bfseries, #5!90!black] at ({#1}, {#2*2.2}) {#4};
    }

    % --- 1. 残り3つのビーカーの並び ---
    % ビーカー 1: 水量 t (<= x)
    \drawFinalBeaker{1.2}{0.42}{ビーカー $1$}{$t \le x$}{teal}
    % ビーカー 2: 水量 s (<= x)
    \drawFinalBeaker{3.2}{0.35}{ビーカー $2$}{$s \le x$}{teal}
    % ビーカー X: 水量 x (< 1/3) (未増量)
    \drawFinalBeaker{5.2}{0.55}{ビーカー $X$}{$x$}{red}

    % --- 2. 上限線 y = x (破線) ---
    \draw[dashed, red!80!black, thick] (0.3, {0.55*2.2}) -- (6.1, {0.55*2.2})
    node[right=2pt, font=\scriptsize, red!90!black] {水面の上限 $x$};

  \end{tikzpicture}
  \caption{ビーカーが残り3個になったときの水量 $t, s, x$ と背理法の矛盾}
  \label{fig:three_beakers_contradiction}
\end{figure}

題意を背理法で示す．
命題の否定「$X$ が最後まで残って，かつ水が増えていない」が成立すると仮定する．
操作を続けてビーカーが残り3個になった時$X$ は残っており，かつ (a) (b) の過程で選択されないので，残り2つのビーカーには $x$ リットル以下の水しか入っていない．これらを$t,s$ [L]とすると，
\begin{align}
  S = t+s+x \le 3x
\end{align}
である．題意より$x<\dfrac{1}{3}$だから\cref{eq:1}も合わせて
\begin{align}
  S \le 3x < 1 \quad \therefore 1 < 1
\end{align}
となり矛盾する．よって背理法により題意は示された．

\paragraph{(2)}

題意を背理法で示す．
命題の否定「$X$ は (a)(b) の過程で選択される」が成立すると仮定する．
(a)(b) の過程では，3個以上のビーカーの中から，入っている水の量が少ない順に2個のビーカーが選ばれる．よって，$X$ が選ばれる時，$x$ [L] 以上の水が入っているビーカーが $X$ 以外に1個以上存在する．
このようなビーカーが $X$ 以外に2個以上存在すると仮定すると
\begin{align}
  S \ge 3x > \dfrac{3}{5}
\end{align}
となり，\cref{eq:1}の $S=1$ に矛盾する．

したがって，$x$ [L] 以上の水が入っているビーカーは $X$ 以外に1個である．それを $Y$ とする．以上の議論より $X$ が選ばれる時はビーカーは残り3個であり，その水の量は $y, z, x$ として，
\begin{align}
  \begin{cases}
    z \le x \le y \\
    S=x+y+z
  \end{cases}\label{eq:2}
\end{align}
と表せる．

ここで，はじめの状態では$X$が最も水の量が多いから，ビーカー$Y$は最初は存在しておらず，$z$ [L]以下の水の入った2個のビーカーの水をあわせて作られる．($z$ [L]より多くの水の入ったビーカーが用いられているならば，$z$ リットルのビーカーは残されない．)
よって，
\begin{align}
  y \le 2z
\end{align}
だから，\cref{eq:2}より
\begin{align}
  S = x + y + z \ge x + \dfrac{3}{2}y \ge x+\dfrac{3}{2}x = \dfrac{5}{2}x
\end{align}
だが，題意より$x>\dfrac{2}{5}$だから
\begin{align}
  S \ge \dfrac{5}{2}x > 1
\end{align}
となり\cref{eq:1}の$S=1$ に矛盾する．

よって背理法によって題意は示された．

{\bf [別解]}

\paragraph{(2)}

$x,y,z$のビーカーがあるところからは以下のような議論も可能である．

$Y$ [L] のビーカーは最初からは存在しないので，このビーカーがつくられる時を考える．この時，$t$個のビーカーがあるとし，水量の少ない順に $a_1, a_2, a_3,\cdots,a_t$ とする．又 $t \ge 4$ としても良い．この時，(a)で $a_1$が，(b)で $a_2$ が各々選ばれるとして良い．
$a_k$ に入っている水量を $b_k$ とすると，$b_1+b_2 = y$ が$x$以上だから
\begin{align}
  b_1+b_2 &\ge x \label{eq:3}
\end{align}
であり，$a_3 \sim a_5$ に入っている水量和は $1-2x$ [L] 以下である．ここから $X$ のビーカーをのぞいたビーカーの内水量の和は $1-2x$ [L] 以下だから，$b_1 \le 1-2x, b_2 \le 1-2x$ が必要で
\begin{align}
  b_1+b_2 &\le 2-4x \label{eq:4}
\end{align}
である．
\cref{eq:1,eq:2}から
\begin{align}
  x \le 2-4x \quad \therefore x \le \dfrac{2}{5}
\end{align}
である．これは $x > \dfrac{2}{5}$ に反し矛盾するから背理法により題意は示された．

\bigskip

\end{document}
